% GENERATED from content/data/enrichment.json by mxbuild.enrichment; do not edit.
\mxunit{enrichment-modular-solutions}{Congruences et derniers chiffres}
\mxsubtitle{Lycée · arithmétique \textperiodcentered{} Solutions détaillées \textperiodcentered{} 03}

\begin{exercise}
Trouver le reste de $-37$ modulo 8.
\end{exercise}
\begin{solution}
$-37 = 8 \times  (-5) + 3$, donc le reste est 3.
\end{solution}

\begin{exercise}
Déterminer le dernier chiffre de $7^{2026}$.
\end{exercise}
\begin{solution}
Le cycle est $7,9,3,1$ et $2026 \equiv  2$ (mod 4). Le dernier chiffre est 9.
\end{solution}

\begin{exercise}
Prouver qu’un carré entier n’est jamais congru à 2 modulo 3.
\end{exercise}
\begin{solution}
Un entier est congru à 0, 1 ou 2. Les carrés correspondants sont 0, 1 et $4 \equiv  1$ modulo 3.
\end{solution}

\begin{exercise}
Résoudre $3x \equiv  1$ (mod 7).
\end{exercise}
\begin{solution}
$3 \times  5 = 15 \equiv  1$. Multiplier par 5 donne $x \equiv  5$ (mod 7). Réciproquement, toutes ces valeurs conviennent.
\end{solution}

\begin{problem}[Chercher, justifier, conclure]
Un entier N est congru à 2 modulo 3 et à 3 modulo 5. Écrire $N = 3k + 2$, déterminer k modulo 5, puis donner toutes les solutions modulo 15. Parmi les solutions entre 0 et 100 inclus, lesquelles sont paires ?
\end{problem}
\begin{solution}
$3k + 2 \equiv  3$ (mod 5), donc $3k \equiv  1$ et $k \equiv  2$ (mod 5), car $3 \times  2 \equiv  1$. Ainsi $N = 15t + 8$. Dans $\intff{0}{100}$, on obtient $8,23,38,53,68,83,98$. Les valeurs paires sont $8,38,68,98$. La substitution vérifie les deux congruences pour toute valeur entière de t.
\end{solution}

